Đối với các quỹ đầu tư và tổ chức tài chính, việc ứng dụng mô hình này mang lại những công cụ phòng vệ mang tính thực tiễn cao:
\begin{itemize}
    \item Thiết lập hạn mức rủi ro (Risk Limits) chính xác hơn: Thay vì dựa vào giả định phân phối chuẩn dễ dẫn đến việc cấp vốn quá tay cho các tài sản tiềm ẩn rủi ro sụp đổ, các tổ chức có thể dùng $CVaR$ dựa trên Copula để tính toán chính xác mức vốn kinh tế (Economic Capital) cần dự trữ.
    \item Kiểm tra sức chịu đựng (Stress Testing): Việc linh hoạt thay đổi các tham số trong mô phỏng Copula cho phép các nhà quản trị tạo ra các kịch bản khủng hoảng cộng hưởng để kiểm tra sức chống chịu của danh mục, từ đó có các biện pháp phòng vệ kịp thời.
\end{itemize}